% =====================================================================
%  U17 —— §10.5 从局部到整体 Langlands 对应
%  覆盖 PDF p.375–p.388（原书页面上的 folio 为 p.365–p.378）
%  底本：E. Frenkel, Langlands Correspondence for Loop Groups, CUP 2007
%  注意：本片段不写 \chapter，接在 ch10-3 之后；开头数段是 §10.4
%        末尾的续文（ch10-3 止于 p.364 页末的半个句子），随后才开始 §10.5。
% =====================================================================

% --------------------------------------------------------------- p.365
获得了有用的解释：它正是范畴 $\hat{\fsl}_2\text{-mod}^{I^0}_\chi$ 的两个不可约对象之扩张空间的射影化。

按照我们的猜想，导出范畴 $D^b(\hat{\fsl}_2\text{-mod}^{I^0}_\chi)$ 等价于 $0$ 处 DG Springer 纤维上拟凝聚层的导出范畴，而该纤维就是 $\Pi T\PP^1$。上面讨论的 Wakimoto 模对应摩天层。那么，拟凝聚层范畴中哪些对象对应不可约 $\hat{\fsl}_2$-模 $\mathbb{V}_0(\chi)$ 与 $\mathbb{V}_{-2}(\chi)$？答案是：$\mathbb{V}_0(\chi)$ 对应 $\PP^1$ 上的平凡线丛 $\cO$，而 $\mathbb{V}_{-2}(\chi)$ 对应复形 $\cO(-1)[1]$，也就是把线丛 $\cO(-1)$ 放在上同调次数 $-1$ 处。$\PP^1$ 上拟凝聚层范畴中的正合列
\[
    0\to \cO(-1)\to \cO\to \cO_p\to 0,
\]
其中 $p\in\PP^1$、$\cO_p$ 是支撑在 $p$ 处的摩天层，变成了范畴 $\hat{\fsl}_2\text{-mod}^{I^0}_\chi$ 中的正合列
\[
    0\to \mathbb{V}_0(\chi)\to W_{u(t)}\to \mathbb{V}_{-2}(\chi)\to 0,
\]
其中 $u(t)$ 是对应于点 $p$ 的方程 (10.4-19) 的解。

现在我们注意到，在范畴 $\hat{\fsl}_2\text{-mod}^{I^0}_\chi$ 中，对象 $\mathbb{V}_0(\chi)$ 与 $\mathbb{V}_{-2}(\chi)$ 之间有一种完全的对称性，因为它们由一个 $\hat{\fsl}_2$ 的对合相联系，而该对合给出 $\hat{\fsl}_2\text{-mod}^{I^0}_\chi$ 的一个自等价。特别地，我们有 $\operatorname{Ext}^1(\mathbb{V}_0(\chi),\mathbb{V}_{-2}(\chi))\simeq\C^2$。但在范畴 $D^b(\operatorname{QCoh}(\PP^1))$ 中却有
\[
    \operatorname{Ext}^1(\cO,\cO(-1)[1]) = H^2(\PP^1,\cO(-1)) = 0 .
\]
这就是为什么范畴 $D^b(\operatorname{QCoh}(\PP^1))$ 不可能等价于 $\hat{\fsl}_2\text{-mod}^{I^0}_\chi$。必须把它换成 $D^b(\operatorname{QCoh}(\Pi T\PP^1))$，即在 DG Springer 纤维 $\Pi T\PP^1$ 上的拟凝聚层范畴。这很好地说明了：在等价 (10.4-14) 中，使用 DG Springer 纤维是必要的。\transnote{$\Pi T\PP^1$ 是“移位切丛”：$\PP^1$ 上的辛形式把切丛与余切丛等同，故 $\Pi T\PP^1\simeq\Pi T^*\PP^1$；其中 $T^*\PP^1$ 是 $\mathfrak{sl}_2$ 幂零锥的 Springer 消解（最小消解），取移位是为了让非约化 DG Springer 纤维的导出范畴与 $\hat{\fsl}_2\text{-mod}^{I^0}_\chi$ 相符。在 $D^b(\operatorname{QCoh}(\PP^1))$ 中 $\operatorname{Ext}^1(\cO,\cO(-1)[1])=H^2(\PP^1,\cO(-1))=0$，而 $\operatorname{Ext}^1(\mathbb{V}_0(\chi),\mathbb{V}_{-2}(\chi))\simeq\C^2$，故拟凝聚层必须取在 DG Springer 纤维上。[来源：原书 p.365；原书 (10.4-14) 前的讨论；§10.4。]}

% [Tutor Note: 这里的关键点是「导出范畴记得扩张」。$\operatorname{Ext}^1(\mathbb{V}_0,\mathbb{V}_{-2})\cong\C^2$ 编码了 Wakimoto 模作为非平凡扩张的信息，而在 $D^b(\operatorname{QCoh}(\PP^1))$ 中两个对象之间没有这种扩张（$H^2(\PP^1,\cO(-1))=0$）。把底空间换成 $\Pi T\PP^1$（即移位切丛 $\Pi T\PP^1$；因 $\PP^1$ 上辛形式把切丛与余切丛等同，它也就是移位余切丛 $\Pi T^*\PP^1$，即 $\mathfrak{sl}_2$ 幂零锥的最小消解）后，额外的一维就出现了。这正是 §10.4 反复出现的主题：局部 Langlands 范畴必须用 DG 概形而非常概形来实现。]

% --------------------------------------------------------------- p.365
\section{从局部到整体 Langlands 对应}

\begin{keypoints}
    \item \textbf{核心思想：}局部的临界水平 $\gghat_{\kappa_c}$-模范畴可以经局部化函子变成曲线 $X$ 上整体的 Hecke 本征层；因此局部的表示论结果能决定整体的几何 Langlands 范畴。
    \item \textbf{数学工具：}局部化函子 $\Delta$（Beilinson--Bernstein、Beilinson--Drinfeld 的构造）、球面与仿射 Hecke 代数、opers、DG Springer 纤维、拟凝聚层。
    \item \textbf{关键结论：}$^L G$-局部系统 $E=(\cF,\nabla)$ 决定 $\Bun_G$ 上的 Hecke 本征层 $\Aut_E$（定理 10.5.1）；带驯顺分歧时局部化把局部范畴送到 $\Bun_{G,(y_i)}$ 上的 Hecke 本征层（定理 10.5.2）。
    \item \textbf{猜想：}(10.5-11)--(10.5-14) 把整体 Langlands 范畴的导出范畴等同于其留数处 DG Springer 纤维上拟凝聚层的导出范畴。
    \item \textbf{记号提醒：}$\chi=(\cF,\nabla,\cF_{LB})\in\Op_{{}^L G}(X)$ 表示 oper；$\Bun_{G,(y_i)}$ 表示带抛物结构的模叠；$\cB_M$ 表示包含 $M$ 的 Borel 子群簇。
\end{keypoints}

我们现在讨论本章所研究的局部 Langlands 对应，对整体几何 Langlands 对应意味着什么。

整体对应的框架我们在 §1.1.6 中已经概述过。回忆一下：在经典图景中，起点是 $\F_q$ 上的一条光滑射影曲线 $X$。记 $F=\F_q(X)$ 为 $X$ 上有理函数构成的域。对 $X$ 的任一闭点 $x$，记 $F_x$ 为 $F$ 在 $x$ 处的完备化，$\cO_x$ 为它的整数环。于是，$X$ 的每一点都附带一个局部域 $F_x$。$F$ 的\textbf{阿代尔}环 $\mathbb{A}=\mathbb{A}_F$ 按定义是各域 $F_x$ 的限制乘积，其中 $x$ 遍历 $X$ 的全部闭点构成的集合 $|X|$（“限制”一词的含义见 §1.1.6）。它含有一个子环 $\mathbb{O}=\prod_{x\in X}\cO_x$。

% --------------------------------------------------------------- p.366
设 $G$ 是 $\F_q$ 上的（分裂）约化群。则对每个 $x\in X$ 有群 $G(F_x)$；它们对一切 $x\in X$ 的限制乘积就是群 $G(\mathbb{A}_F)$。由于 $F\subset\mathbb{A}_F$，群 $G(F)$ 是 $G(\mathbb{A}_F)$ 的子群。

整体 Langlands 对应的一侧是满足某些性质的同态 $\sigma:W_F\to{}^L G$（或者还带有某些更精细的数据，如 [Arthur 1989] 中那样）。我们期望能对每个 $\sigma$ 附上一个 $GL_n(\mathbb{A}_F)$ 的\textbf{自守表示} $\pi$。\footnote{本节中，我们按惯例滥用记号，用同一个符号表示群的表示与它底层的向量空间。}“自守”一词大致是说：该表示可以实现在商空间 $GL_n(F)\backslash GL_n(\mathbb{A})$ 上的某个合理的函数空间里（$GL_n(\mathbb{A})$ 从右侧作用）。我们不打算把它讲精确。一般地，我们期望附给 $\sigma$ 的不是一个而是若干个自守表示，它们是前面讨论过的 L-包的整体类似物（见 [Arthur 1989]）。另一个复杂之处是：给定的不可约自守表示在 $GL_n(F)\backslash GL_n(\mathbb{A})$ 上函数空间中的重数可能大于 $1$。这里我们基本上忽略所有这些麻烦，因为我们的主要兴趣在几何理论（注意：当 $G=GL_n$ 时这些问题都不出现）。

任一不可约自守表示总可以分解为限制张量积 $\bigotimes'_{x\in X}\pi_x$，其中每个 $\pi_x$ 是 $G(F_x)$ 的不可约表示。此外，对除有限多个之外的一切 $x\in X$，因子 $\pi_x$ 是 $G(F_x)$ 的\textbf{无分歧表示}：它含有在极大紧子群 $K_{0,x}=G(\cO_x)$ 下不变的非零向量（见 §10.3.1）。取这样一个向量 $v_x\in\pi_x$（它在相差一个标量的意义下唯一）。所谓“限制”，指的是我们考虑形如 $\otimes_{x\in X}u_x$ 的向量的张成，其中 $u_x\in\pi_x$，且对除有限多个之外的一切 $x\in X$ 有 $u_x=v_x$。

整体 Langlands 对应的一个重要性质是它与局部对应相容。我们可以把每个局部域 $F_x$ 的 Weil 群 $W_{F_x}$ 嵌入整体 Weil 群 $W_F$。这样的嵌入不唯一，但它在 $W_F$ 中共轭的意义下良定义。因此，同态 $\sigma:W_F\to{}^L G$ 的等价类给出同态 $\sigma_x:W_{F_x}\to{}^L G$ 的一个良定义的等价类。我们将对 $\sigma$ 加上条件：对除有限多个之外的一切 $x\in X$，同态 $\sigma_x$ 是无分歧的（见 §10.3.1）。

由局部 Langlands 对应，可以对 $\sigma_x$ 附上 $G(F_x)$ 的不可约光滑表示的等价类 $\pi_x$。\footnote{这里我们考虑的是从 Weil 群 $W_{F_x}$ 到 ${}^L G$ 的 $\ell$-adic 同态，因此不需要从 Weil 群过渡到 Weil--Deligne 群。}此外，无分歧的 $\sigma_x$ 将对应无分歧的不可约表示 $\pi_x$。局部对应与整体对应之间的相容性，说的就是：对应 $\sigma$ 的 $G(\mathbb{A})$ 的自守表示应当同构于限制张量积 $\bigotimes'_{x\in X}\pi_x$。用图式表示，就是：

% --------------------------------------------------------------- p.367
\[
\begin{array}{c}
\sigma \;\overset{\text{整体}}{\longleftrightarrow}\;
\pi=\bigotimes{}'_{x\in X}\pi_x\\[5pt]
\sigma_x \;\overset{\text{局部}}{\longleftrightarrow}\;
\pi_x .
\end{array}
\]

本节我们讨论这一“局部到整体”原理在几何图景中的类似物，以及我们的局部结果与猜想对整体几何 Langlands 对应的推论。我们特别关注无分歧与驯顺分歧的 Langlands 参数。

\subsection{无分歧情形}

重要的特殊情形是 $\sigma:W_F\to{}^L G$ 处处无分歧。此时对每个 $x\in X$，相应的同态 $\sigma_x:W_{F_x}\to{}^L G$ 无分歧，因而如 §10.3.1 所说明，它对应 ${}^L G$ 中的一个半单共轭类 $\gamma_x$，即 Frobenius 元在 $\sigma_x$ 下的像。这个共轭类又给出 $G(F_x)$ 的一个无分歧不可约表示 $\pi_x$，其中存在唯一的（相差一个标量）向量 $v_x$，满足 $G(\cO_x)v_x=v_x$。球面 Hecke 代数 $H(G(F_x),G(\cO_x))\simeq\Rep{}^L G$ 按公式 (10.3-5) 作用在这个向量上：
\begin{equation}
    H_{V,x}\star v_x = \Tr(\gamma_x,V)\,v_x,
    \qquad [V]\in\Rep{}^L G.
    \tag{10.5-1}
\end{equation}
这些向量的张量积 $v=\otimes_{x\in X}v_x$ 是 $\pi=\bigotimes'_{x\in X}\pi_x$ 中的一个 $G(\mathbb{O})$-不变向量；按整体 Langlands 猜想，它是自守的。这意味着 $\pi$ 实现在 $G(F)\backslash G(\mathbb{A}_F)$ 上的函数空间中。在这个实现里，向量 $v$ 对应 $G(F)\backslash G(\mathbb{A}_F)$ 上的一个右 $G(\mathbb{O})$-不变函数，等价地说，对应双商
\begin{equation}
    G(F)\backslash G(\mathbb{A}_F)/G(\mathbb{O})
    \tag{10.5-2}
\end{equation}
上的一个函数。

% [Tutor Note: 公式 (10.5-1) 是 Satake 同构下「Hecke 本征向量」的定义式：$H_{V,x}$ 是球面 Hecke 代数中对应表示 $V$ 的元素，在 Satake 同构 $H(G(F_x),G(\cO_x))\simeq\Rep{}^L G$ 下 $\Rep{}^L G$ 的表示环与 Hecke 代数同构；$\gamma_x$ 是未分歧参数对应的半单共轭类，$\Tr(\gamma_x,V)$ 就是通常的 Hecke 本征值。注意 $\star$ 表示 Hecke 代数的卷积作用，不是普通的数乘；把 $\gamma_x$ 与 $\sigma_x(\mathrm{Frob})$ 的共轭类等同是 Langlands 对应的核心约定，核对时不要漏掉 Frobenius 的几何/算术归一化差异。]

于是，无分歧的整体 Langlands 参数 $\sigma$ 给出 (10.5-2) 上的一个函数。这个函数就是对应 $\sigma$ 的\textbf{自守函数}，我们记作 $f_\pi$。由于它对应 $G(\mathbb{A}_F)$ 的不可约表示 $\pi$ 中的一个向量，整个表示 $\pi$ 都可以由这个函数重构。因此，从 $\pi$ 过渡到 $f_\pi$ 不丢失任何信息。

由于 $v\in\pi$ 是 Hecke 算子的本征向量，按公式 (10.5-1) 我们得到：函数 $f_\pi$ 是双商 (10.5-2) 上的\textbf{Hecke 本征函数}。事实上，局部 Hecke 代数 $H(G(F_x),G(\cO_x))$ 自然地（从右侧）作用在 (10.5-2) 上的函数空间上，而 $f_\pi$ 是这个作用的本征函数，它满足同样的性质 (10.5-1)。\transnote{这里的“本征值”必须是 $X$ 上的一个局部系统：几何上 Hecke 算子的本征值取平坦 ${}^L G$-丛 $E=(\cF,\nabla)$（局部自由层带平坦联络），而局部 Hecke 本征值要求对曲线 $X$ 的每个点 $x$ 同时指定本征值并满足 (10.5-1)。这些局部本征值必须来自同一个整体参数（经典情形为同一个 $\sigma:W_F\to{}^L G$，几何情形为同一个局部系统 $E$）——这就是整体本征条件中的“相容性”。只在单点取本征条件得到的是局部对象；对一切 $x\in X$ 同时取，才能整体地刻画 $\pi$（对应 $f_\pi$）。[来源：原书 p.367；§10.3.1；(10.5-1)。]}

% --------------------------------------------------------------- p.368
总结一下：经典图景中的无分歧整体 Langlands 对应，可以看作无分歧同态 $\sigma:W_F\to{}^L G$ 与 (10.5-2) 上的 Hecke 本征函数之间的对应（为使它更精确，还需对 $\sigma$ 加上某种不可约性条件，但我们忽略它）。

当 $X$ 是复代数曲线时，这个对应的几何类似物应当是什么？

正如 §1.2.1 所说明的，无分歧同态 $W_F\to{}^L G$ 的几何类似物是同态 $\pi_1(X)\to{}^L G$；或者等价地说（因假定 $X$ 紧），是 $X$ 上带全纯联络的全纯 ${}^L G$-丛（它自动给出平坦联络，见 §1.2.3）。因此，整体几何 Langlands 对应应当把 $X$ 上的平坦全纯 ${}^L G$-丛对应到双商 (10.5-2) 的某个几何版本上的几何对象。正如 §1.3.2 所论证的，这应当是点集为 (10.5-2) 的某个代数簇上的 D-模。

现在已知：(10.5-2) 与 $X$ 上 $G$-丛的同构类集合一一对应。这个关键结果属于 A. Weil（例如见 [2007]）。这提示 (10.5-2) 是 $X$ 上 $G$-丛模空间的点集。遗憾的是，它一般不是代数簇，而是代数叠：局部地看，它是一个代数簇用一个代数群的作用取商。我们把它记作 $\Bun_G$。在 [Beilinson and Drinfeld 1997] 中，D-模理论已经在 $\Bun_G$ 这样的代数叠的框架下发展起来，因此我们可以把它用于我们的目的。于是，我们希望把 $X$ 上的平坦全纯 ${}^L G$-丛 $E$ 对应到 $\Bun_G$ 上的一个 D-模 $\Aut_E$。这个 D-模应当满足 Hecke 本征函数条件的类似物，从而成为以 $E$ 为本征值的\textbf{Hecke 本征层}。这一概念在 [Frenkel 2007] 中（沿袭 [Beilinson and Drinfeld 1997]）有详细说明，读者可参阅其中的细节。

于是，无分歧的整体几何 Langlands 对应可以概括为下面的图式：
\[
\begin{array}{ccc}
    \boxed{\ \text{平坦 ${}^L G$-丛（$X$ 上）}\ } & \longrightarrow &
    \boxed{\ \text{$\Bun_G$ 上的 Hecke 本征层}\ }\\[7pt]
    E & \mapsto & \Aut_E
\end{array}
\]
（这里应当对 $E$ 加上某种不可约性条件）。

无分歧的整体几何 Langlands 对应已在 [Frenkel et al. 2002; Gaitsgory 2004] 中对 $G=GL_n$ 及任意不可约 $GL_n$-局部系统 $E$ 得到证明；在 [Beilinson and Drinfeld 1997] 中对任意单李群 $G$ 以及那些具有 oper 结构的 ${}^L G$-局部系统得到证明。本节我们讨论第二条途径。

为说明其动机，我们提出如下问题：

% --------------------------------------------------------------- p.369
\textbf{如何把这个整体对应与上面讨论的局部几何 Langlands 对应联系起来？}

回答这个问题的关键要素，是一个从 $\gghat$-模到 $\Bun_G$ 上（扭变）D-模的\textbf{局部化函子}。为构造这个函子，注意对单李群 $G$，模叠 $\Bun_G$ 还有另一种双商实现。具体地说，取 $X$ 上的点 $x$，记 $\cO_x$ 为完备局部环、$\cK_x$ 为其分式域。令 $G(\cK_x)$ 为围绕 $x$ 的去心圆盘 $D^\times_x=\Spec\cK_x$ 对应的形式圈群，$G_{\mathrm{out}}$ 为代数映射 $X\setminus x\to G$ 构成的群，它自然地是 $G(\cK_x)$ 的子群。则 $\Bun_G$ 同构于双商
\begin{equation}
    \Bun_G \simeq G_{\mathrm{out}}\backslash G(\cK_x)/G(\cO_x).
    \tag{10.5-3}
\end{equation}
这是 (10.5-2) 的“单点”版本。在 $\C$-点的层面上这并不难证明，但这个同构在代数叠的层面上同样成立（见 [Beauville and Laszlo 1995; Drinfeld and Simpson 1995]）。

我们需要的局部化函子，是下述一般构造的特殊情形。设 $(\fg,K)$ 是 Harish--Chandra 对，$\fg\text{-mod}^K$ 是 $K$-等变 $\fg$-模范畴（见 §10.2.2）。对子群 $H\subset G$，记 $\cD_{H\backslash G/K}\text{-mod}$ 为 $H\backslash G/K$ 上 D-模构成的范畴。则存在局部化函子 [Beilinson and Bernstein 1993; Beilinson and Drinfeld 1997]（另见 [Frenkel 2007; Frenkel and Ben-Zvi 2004]）\transnote{局部化函子的一般机制来自 Beilinson--Bernstein：把 $K$-等变 $\fg$-模实现为 $H\backslash G/K$ 上的（扭变）$D$-模。这里的关键是把模叠 $\Bun_G$ 写成 (10.5-3) 式的双商 $G_{\mathrm{out}}\backslash G(\cK_x)/G(\cO_x)$，于是同一形式体系可以把 $\gghat_{\kappa_c,x}$-模局部化为 $\Bun_G$ 上的 $D$-模；中心扩张指数化到群后出现的 $\C^\times$-丛给出线丛 $\cL$，其临界水平版本正是 $K^{1/2}$。[来源：原书 p.369；[Beilinson and Bernstein 1993]；[Beilinson and Drinfeld 1997]。]}
\[
    \Delta:\ \fg\text{-mod}^K \to \cD_{H\backslash G/K}\text{-mod}.
\]
现在设 $\gghat$ 是 $\fg$ 的一维中心扩张，且它在限制到李子代数 $\Lie K$ 与 $\Lie H$ 上时变成平凡扩张。假设这个中心扩张可以提升为相应李群 $G$ 的中心扩张 $\hat{G}$。于是我们得到 $H\backslash G/K$ 上的一个 $\C^\times$-丛 $H\backslash\hat{G}/K$。令 $\cL$ 为相应的线丛，$\cD_\cL$ 为作用在 $\cL$ 上的微分算子层。则我们有一个函子
\[
    \Delta_\cL:\ \gghat\text{-mod}^K \to \cD_\cL\text{-mod}.
\]

在我们的情形中，取李代数 $\gghat_{\kappa_c,x}$，即圈代数 $\fg\otimes\cK_x$ 的临界中心扩张，并取 $G(\cK_x)$ 的子群 $K=G(\cO_x)$ 与 $H=G_{\mathrm{out}}$。已知（例如见 [Beilinson and Drinfeld 1997]）相应的线丛 $\cL$ 是 $\Bun_G$ 上典范线丛的平方根 $K^{1/2}$。\footnote{回忆按我们的假定 $G$ 是单连通的；此时平方根唯一。}我们把 $\Bun_G$ 上相应的扭变微分算子层记作 $\cD_{\kappa_c}$。于是

% --------------------------------------------------------------- p.370
我们得到局部化函子
\[
    \Delta_{\kappa_c,x}:\ \gghat_{\kappa_c,x}\text{-mod}^{G(\cO_x)}
    \to \cD_{\kappa_c}\text{-mod}.
\]
我们再把它限制到范畴 $\gghat_{\kappa_c,x}\text{-mod}^{G(\cO_x)}_{\chi_x}$ 上，其中
\[
    \chi_x\in\Op_{{}^L G}(D^\times_x).
\]
由引理 10.3.2，这个范畴非零，仅当 $\chi_x\in\Op^{\lambda}_{{}^L G}(D_x)$ 对某个 $\lambda\in P^+$ 成立。为简单起见，设 $\lambda=0$，从而 $\chi_x\in\Op_{{}^L G}(D_x)$。此时由定理 10.3.3，范畴 $\gghat_{\kappa_c,x}\text{-mod}^{G(\cO_x)}_{\chi_x}$ 等价于向量空间范畴。它在同构意义下唯一的不可约对象，就是真空模
\[
    \mathbb{V}_{0,x} = \Ind^{\gghat_{\kappa_c,x}}_{\fg(\cO_x)\oplus\C\mathbf{1}}\C
\]
对相应于 $\chi_x$ 的中心特征标取商所得的对象 $\mathbb{V}_0(\chi_x)$。因此只需确定 $\Delta_{\kappa_c,x}(\mathbb{V}_0(\chi_x))$。

下述定理属于 Beilinson 与 Drinfeld [1997]。

\begin{theorem}
    (1) $\cD_{\kappa_c}$-模 $\Delta_{\kappa_c,x}(\mathbb{V}_0(\chi_x))$ 非零，当且仅当存在 $X$ 上的整体 ${}^L\fg$-oper\transnote{原书此处印作“整体 ${}^L\mathfrak{g}$-oper”，按体例应为“整体 ${}^L G$-oper”：oper 是按群 ${}^L G$ 定义的，同句后半正作 $\chi\in\Op_{{}^L G}(X)$，第 4 章与 p.371 亦一律用 ${}^L G$。译文保留原书印法，未擅改。[来源：原书 p.370（PDF p.380）页图；原书 p.371；第 4 章。]}，$\chi\in\Op_{{}^L G}(X)$，使得 $\chi_x\in\Op_{{}^L G}(D_x)$ 是 $\chi$ 到 $D_x$ 的限制。

    (2) 若上述条件成立，则 $\Delta_{\kappa_c,x}(\mathbb{V}_0(\chi_x))$ 只依赖于 $\chi$，并在下述意义下与 $x$ 的选取无关：对任意另一点 $y\in X$，若 $\chi_y=\chi|_{D_y}$，则 $\Delta_{\kappa_c,x}(\mathbb{V}_0(\chi_x))\simeq\Delta_{\kappa_c,y}(\mathbb{V}_0(\chi_y))$。

    (3) 对任意 $\chi=(\cF,\nabla,\cF_{LB})\in\Op_{{}^L G}(X)$，$\cD_{\kappa_c}$-模 $\Delta_{\kappa_c,x}(\mathbb{V}_0(\chi_x))$ 是以 $E_\chi=(\cF,\nabla)$ 为本征值的非零 Hecke 本征层。
\end{theorem}

% [Tutor Note: 定理 10.5.1 (1) 中的「整体 ${}^L\fg$-oper」与紧随其后的 $\chi\in\Op_{{}^L G}(X)$ 不匹配：oper 是按群 ${}^L G$（而非李代数 ${}^L\fg$）定义的，此处疑为「${}^L G$-oper」的印刷之误。原样保留，未擅改。另注意结论 (2) 是「与基点无关」这一整体性陈述，它保证局部化给出的对象真正是 $X$ 上的层。]

于是，对任意 $\chi\in\Op_{{}^L G}(X)$，$\cD_{\kappa_c}$-模 $\Delta_{\kappa_c,x}(\mathbb{V}_0(\chi_x))$ 就是在整体几何 Langlands 对应 10.5.1 之下对应 ${}^L G$-局部系统 $E_\chi$ 的那个所求的 Hecke 本征层 $\Aut_{E_\chi}$。\footnote{更确切地说，$\Aut_{E_\chi}$ 是 D-模 $\Delta_{\kappa_c,x}(\mathbb{V}_0(\chi_x))\otimes K^{-1/2}$，但此处及下文我们都忽略 $K^{1/2}$ 的扭变。}关于这个定理（出自 [Beilinson and Drinfeld 1997]）的证明要点，见 [Frenkel 2007]。关键之处在于：对一切点 $y\in X$，相应于 $y$ 的不可约模 $\mathbb{V}_0(\chi_y)$（其中 $\chi_y=\chi|_{D_y}$）都是 Hecke 本征模，正如我们在 §10.3.5 所见。\transnote{定理 10.5.1 的整体意义在于：对每个点 $x$，局部化只给出 $X$ 上同一个 Hecke 本征层的局部片段，而 (2) 的“与 $x$ 的选取无关”正是把这些片段拼成整体层的相容性条件——换言之，$\chi$ 在每个点 $y$ 处的不可约模 $\mathbb{V}_0(\chi_y)$ 必须同时是本征模，才可能来自同一个整体对象。另按脚注 15，本征层 $\Aut_{E_\chi}$ 与 $\Delta_{\kappa_c,x}(\mathbb{V}_0(\chi_x))$ 相差 $K^{-1/2}$ 的扭变，书中约定忽略该扭变。[来源：原书 p.370 及脚注 15；§10.3.5。]}

这个构造的一个缺点是：并非 $X$ 上所有 ${}^L G$-局部系统都具有 oper 结构。事实上，在我们假定 ${}^L G$ 是伴随型群的条件下，光滑射影曲线 $X$ 上具有 oper 结构的 ${}^L G$-局部系统（即平坦丛 $(\cF,\nabla)$）对应 $X$ 上唯一的 ${}^L G$-丛 $\cF_{\mathrm{oper}}$（见 §4.2.4）。此外，对每个这样的平坦丛 $(\cF_{\mathrm{oper}},\nabla)$，存在

% --------------------------------------------------------------- p.371
唯一的 ${}^L B$-约化 $\cF_{\mathrm{oper},LB}$ 满足 oper 条件。因此，$\Op_G(D)$ 是遗忘映射
\[
    \Loc_{{}^L G}(X)\to\Bun_{{}^L G},
    \qquad (\cF,\nabla)\mapsto \cF
\]
在 oper 丛 $\cF_{\mathrm{oper}}$ 之上的纤维。\transnote{此处原书印作 $\Op_G(D)$，按上下文应为 $\Op_{{}^L G}(X)$：这是遗忘映射 $\Loc_{{}^L G}(X)\to\Bun_{{}^L G}$ 在 oper 丛 $\cF_{\mathrm{oper}}$ 之上的纤维，而 oper 丛属于 ${}^L G$（第 4 章与 p.370 均用 $\Op_{{}^L G}(X)$）。译文保留原书印法，未擅改。[来源：原书 p.371（PDF p.381）页图；原书 p.370；第 4 章。]}

% [Tutor Note: 此处原文印作 $\Op_G(D)$，但上下文谈的是 $X$ 上整体的 oper，按记号体系应为 $\Op_{{}^L G}(X)$（下文及第 4 章均用后者）。这是原书的疑似笔误，译文保留原样，未擅改，已登记报告。]

对满足 $\cF\neq\cF_{\mathrm{oper}}$ 的 ${}^L G$-局部系统 $E=(\cF,\nabla)$，上述构造可以如下修改（见 [Frenkel 2007] 中的讨论，基于 Beilinson 与 Drinfeld 未发表的工作）。假设我们能选取一个 ${}^L B$-约化 $\cF_{LB}$，使它在某个有限点集 $y_1,\dots,y_n$ 之外满足 oper 条件，并且相应的 oper $\chi$（定义在 $X\setminus\{y_1,\dots,y_n\}$ 上）到 $D^\times_{y_i}$ 的限制 $\chi_{y_i}$ 属于 $\Op^{\lambda_i}_{{}^L G}(D_{y_i})\subset\Op_{{}^L G}(D^\times_{y_i})$，其中 $\lambda_i\in P^+$。那么，把多点版本的局部化函子作用于 Weyl 模 $\mathbb{V}_{\lambda_i,y_i}$ 的商 $\mathbb{V}_{\lambda_i}(\chi_{y_i})$ 的张量积，就得到对应 $E$ 的 Hecke 本征层（见 [Frenkel 2007]）。

定理 10.5.1 的主要教益是：在几何图景中，局部化函子是强有力的工具，能把局部 Langlands 范畴（如 $\gghat_{\kappa_c,x}\text{-mod}^{G(\cO_x)}_{\chi_x}$）转换为整体的 Hecke 本征层范畴。因此，重点转移到对 $\gghat_{\kappa_c,x}$-模的局部范畴的研究。通过研究局部范畴，我们可以对整体范畴有大量了解。这是一种在经典 Langlands 对应中没有明显类似物的新现象。

在眼下这个情形中，范畴 $\gghat_{\kappa_c,x}\text{-mod}^{G(\cO_x)}_{\chi_x}$ 非常简单：它只有一个不可约对象 $\mathbb{V}_0(\chi_x)$。这就是为什么只需考虑它在局部化函子下的像，而它恰好就是所求的 Hecke 本征层 $\Aut_{E_\chi}$。对于带分歧的一般 oper，相应的局部范畴更为复杂（如上所见），相应的 Hecke 本征层范畴也如此。下一节我们将考虑这些范畴的例子。

\subsection{带驯顺分歧的整体 Langlands 对应}

先看经典图景中带分歧的整体 Langlands 对应。设有同态 $\sigma:W_F\to{}^L G$，它在 $X$ 的有限多个点 $y_1,\dots,y_n$ 处分歧。我们期望这样的 $\sigma$ 对应一个自守表示 $\bigotimes'_{x\in X}\pi_x$（更确切地说，是一个 L-包）。这里 $\pi_x$ 对一切 $x\in X\setminus\{y_1,\dots,y_n\}$ 仍无分歧，但在 $y_1,\dots,y_n$ 处分歧，即 $\pi_{y_i}$ 中的 $G(\cO_{y_i})$-不变向量空间为零。特别地，考虑每个 $\sigma_{y_i}:W_{F_{y_i}}\to{}^L G$ 都\textbf{驯顺分歧}的特殊情形（定义见 §10.4.1）。此时按 §10.4.1 中陈述的结果，$G(F_{y_i})$ 的相应 L-包中含有一个不可约

% --------------------------------------------------------------- p.372
表示 $\pi_{y_i}$，它带有相对于 Iwahori 子群 $I_{y_i}$ 的非零不变向量。对每个点 $y_i$ 取这样一个表示。

考虑子空间
\begin{equation}
    \bigotimes_{i=1}^n \pi_{y_i}^{I_{y_i}}
    \otimes \bigotimes_{x\neq y_i} v_x
    \subset \bigotimes{}'_{x\in X}\pi_x,
    \tag{10.5-4}
\end{equation}
其中 $v_x$ 是 $\pi_x$ 中的 $G(\cO_x)$-向量，$x\neq y_i$、$i=1,\dots,n$。由于 $\bigotimes'_{x\in X}\pi_x$ 实现在 $G(F)\backslash G(\mathbb{A}_F)$ 上的函数空间中，我们得到：子空间 (10.5-4) 实现在双商
\begin{equation}
    G(F)\backslash G(\mathbb{A}_F)\Big/
    \prod_{i=1}^n I_{y_i}\times\prod_{x\neq y_i}G(\cO_x)
    \tag{10.5-5}
\end{equation}
上的函数空间中。

球面 Hecke 代数 $H(G(F_x),G(\cO_x))$（$x\neq y_i$）作用在子空间 (10.5-4) 上，且 (10.5-4) 的所有元素都是这些代数的本征函数（它们满足公式 (10.5-1)）。在点 $y_i$ 处，代替交换的球面 Hecke 代数 $H(G(F_{y_i}),G(\cO_{y_i}))$ 的作用，我们有非交换仿射 Hecke 代数 $H(G(F_{y_i}),I_{y_i})$ 的作用。于是我们得到 (10.5-5) 上函数空间的一个子空间，它由相对于球面 Hecke 代数 $H(G(F_x),G(\cO_x))$（$x\neq y_i$）的 Hecke 本征函数组成，并且实现 $\bigotimes_{i=1}^n H(G(F_{y_i}),I_{y_i})$ 上的一个模（该模不可约，因为每个 $\pi_{y_i}$ 不可约）。\transnote{驯顺分歧处 Iwahori 子群的双重角色：局部 Langlands 参数驯顺分歧时 $\pi_{y_i}$ 的 $I_{y_i}$-不变向量非零（§10.4.1），于是 $y_i$ 处起作用的不再是可交换的球面 Hecke 代数 $H(G(F_{y_i}),G(\cO_{y_i}))$，而是非交换的仿射 Hecke 代数 $H(G(F_{y_i}),I_{y_i})$——Iwahori 球表示正是它的表示论来源。几何上这一步把 $\Bun_{G,(y_i)}$ 上的对象细化为 $\Bun_{G,y}$ 上的 $\cD_{\kappa_c,I_y}$-模，局部化函子 (10.5-6) 才能把局部范畴送到整体范畴。[来源：原书 p.372；§10.4.1；§10.4.5。]}

这个子空间封装自守表示 $\bigotimes'_{x\in X}\pi_x$ 的方式，正如自守函数 $f_\pi$ 封装无分歧自守表示。区别在于：无分歧情形下函数 $f_\pi$ 张成 $\bigotimes'_{x\in X}\pi_x$ 中极大紧子群 $G(\mathbb{O})$ 的不变向量的一维空间；而驯顺分歧情形下子空间 (10.5-4) 一般是多维向量空间。

现在来看这在几何图景中如何体现。如前所述，在点 $y_1,\dots,y_n\in X$ 处驯顺分歧的同态 $\sigma:W_F\to{}^L G$ 的类似物，现在是一个局部系统 $E=(\cF,\nabla)$，其中 $\cF$ 是 $X$ 上的 ${}^L G$-丛，$\nabla$ 是 $X$ 上在 $y_1,\dots,y_n$ 处具有正则奇点、且围绕这些点的单值性为幂幺的联络。我们把这样的局部系统称为在 $y_1,\dots,y_n$ 处\textbf{驯顺分歧}。整体几何 Langlands 对应应当对这种局部系统附上什么？显然，我们需要找到一个几何对象，来替换实现在 (10.5-5) 上函数空间中的有限维向量空间 (10.5-4)。

正如 (10.5-2) 是 $G$-丛模叠 $\Bun_G$ 的点集，双商 (10.5-5) 是 $X$ 上带抛物结构的 $G$-丛模叠 $\Bun_{G,(y_i)}$ 的点集，抛物结构给在 $y_i$（$i=1,\dots,n$）处。按定义，

% --------------------------------------------------------------- p.373
$G$-丛 $\cP$ 在 $y\in X$ 处的一个抛物结构，是把 $\cP$ 在 $y$ 处的纤维 $\cP_y$ 约化到一个 Borel 子群 $B\subset G$。因此如前，我们得到：(10.5-4) 的恰当替换是 $\Bun_{G,(y_i)}$ 上的 D-模范畴。与无分歧情形一样，我们有 $\Bun_{G,(y_i)}$ 上 Hecke 本征层的概念。但由于 Hecke 函子现在是用 $X\setminus\{y_1,\dots,y_n\}$ 上（而不是像以前那样在 $X$ 上）的 Hecke 对应定义的，Hecke 算子的“本征值”现在是 $X\setminus\{y_1,\dots,y_n\}$ 上的 ${}^L G$-局部系统（而不是 $X$ 上的）。于是我们得到：整体几何 Langlands 对应现在应当把 $X$ 上在点 $y_1,\dots,y_n$ 处驯顺分歧的 ${}^L G$-局部系统 $E$，对应到 $\Bun_{G,(y_i)}$ 上以 $E|_{X\setminus\{y_1,\dots,y_n\}}$ 为本征值的 D-模范畴 $\mathcal{A}ut_E$\transnote{务必区分两类记号：无分歧情形得到单个 Hecke 本征层，记正体 $\Aut_E$；带分歧情形得到的是范畴，记花体 $\mathcal{A}ut_E$（下文还有 $\mathcal{A}ut^{I_y}_E$、$\mathcal{A}ut^{I_y,m}_E$）。原因是此时 (10.5-4) 是多维向量空间，几何对象必须从“一个层”升级为“一个范畴”才能与之对应。[来源：原书 p.368、p.373。]}：
\[
    E \mapsto \mathcal{A}ut_E .
\]

% [Tutor Note: 这里出现了两类不同的对象，务必区分：单个的 Hecke 本征层记作 $\Aut_E$（正体），而带分歧情形下得到的是 \emph{范畴} $\mathcal{A}ut_E$（花体）。原因在于此时 (10.5-4) 是多维的，几何上必须用范畴而非单个层来实现。下文的 $\mathcal{A}ut^{I_y}_E$、$\mathcal{A}ut^{I_y,m}_E$ 都是范畴。]

我们将用无分歧情形所用局部化函子的推广来构造这些范畴（见 [Frenkel and Gaitsgory 2006c]）。为记号简单，设 ${}^L G$-局部系统 $E=(\cF,\nabla)$ 只在单点 $y\in X$ 处驯顺分歧。假设 $X\setminus y$ 上的这个局部系统具有 ${}^L G$-oper $\chi=(\cF,\nabla,\cF_{LB})$ 的结构，且它到去心圆盘 $D^\times_y$ 的限制 $\chi_y$ 属于由幂零 ${}^L G$-oper 构成的子空间 $\Op^{\mathrm{nilp}}_{{}^L G}(D_y)$。

对单李群 $G$，模叠 $\Bun_{G,y}$ 有与 (10.5-3) 类似的实现：
\[
    \Bun_{G,y}\simeq G_{\mathrm{out}}\backslash G(\cK_y)/I_y.
\]
令 $\cD_{\kappa_c,I_y}$ 为 $\Bun_{G,y}$ 上作用在临界水平对应线丛上的扭变微分算子层（它是自然投影 $\Bun_{G,y}\to\Bun_G$ 下 $\Bun_G$ 上典范线丛平方根 $K^{1/2}$ 的拉回）。应用上一节的形式体系，我们得到局部化函子
\[
    \Delta_{\kappa_c,I_y}:\ \gghat_{\kappa_c,y}\text{-mod}^{I_y}
    \to \cD_{\kappa_c,I_y}\text{-mod}.
\]
不过，为了与上面所得的结果相衔接，我们还要考虑更大的范畴 $\gghat_{\kappa_c,y}\text{-mod}^{I^0_y}$，即 $I^0_y$-等变模构成的范畴，其中 $I^0_y=[I_y,I_y]$。令
\[
    \Bun'_{G,y}=G_{\mathrm{out}}\backslash G(\cK_y)/I^0_y,
\]
并令 $\cD_{\kappa_c,I^0_y}$ 为 $\Bun'_{G,y}$ 上作用在 $\Bun_G$ 上线丛 $K^{1/2}$ 的拉回上的扭变微分算子层。记 $\cD_{\kappa_c,I^0_y}\text{-mod}$ 为 $\cD_{\kappa_c,I^0_y}$-模范畴。应用一般形式体系，我们得到局部化函子
\begin{equation}
    \Delta_{\kappa_c,I^0_y}:\ \gghat_{\kappa_c,y}\text{-mod}^{I^0_y}
    \to \cD_{\kappa_c,I^0_y}\text{-mod}.
    \tag{10.5-6}
\end{equation}
我们指出，§10.4.5 中讨论的范畴化仿射 Hecke 代数 $\cH(G(\cK_y),I_y)$ 的一个版本自然地作用在上述

% --------------------------------------------------------------- p.374
范畴的导出范畴上，而函子 $\Delta_{\kappa_c,I_y}$ 与 $\Delta_{\kappa_c,I^0_y}$ 交织这些作用。等价地说，这个函子交织 ${}^L G$ 相应的仿射辫群在这两个范畴上的作用（如 [Bezrukavnikov 2006] 中那样）。

现在把函子 $\Delta_{\kappa_c,I_y}$ 与 $\Delta_{\kappa_c,I^0_y}$ 分别限制到子范畴
\[
    \gghat_{\kappa_c,y}\text{-mod}^{I_y}_{\chi_y}
    \qquad\text{和}\qquad
    \gghat_{\kappa_c,y}\text{-mod}^{I^0_y}_{\chi_y}
\]
上。利用与 [Beilinson and Drinfeld 1997] 中相同的论证，我们得到下述与定理 10.5.1 类似的结论。

\begin{theorem}
    固定 $\chi_y\in\Op^{\mathrm{nilp}}_{{}^L G}(D_y)$，设 $M$ 是范畴 $\gghat_{\kappa_c,y}\text{-mod}^{I_y}_{\chi_y}$（相应地，$\gghat_{\kappa_c,y}\text{-mod}^{I^0_y}_{\chi_y}$）中的对象。则：

    (1) $\Delta_{\kappa_c,I_y}(M)=0$（相应地，$\Delta_{\kappa_c,I^0_y}(M)=0$），除非 $\chi_y$ 是 $X\setminus y$ 上某个正则 oper $\chi=(\cF,\nabla,\cF_{LB})$ 到 $D^\times_y$ 的限制。

    (2) 在后一情形下，$\Delta_{\kappa_c,y}(M)$\transnote{原书此处印作 $\Delta_{\kappa_c,y}(M)$，按 (1) 与上下文应为 $\Delta_{\kappa_c,I_y}(M)$：本定理只涉及 $I_y$ 版与 $I^0_y$ 版的两个局部化函子（(10.5-6) 与 p.373 的定义），并不存在以点 $y$ 为下标的 $\Delta_{\kappa_c,y}$。译文保留原书印法，未擅改。[来源：原书 p.374（PDF p.384）页图；原书 p.373。]}（相应地，$\Delta_{\kappa_c,I^0_y}(M)$）是以 $E_\chi=(\cF,\nabla)$ 为本征值的 Hecke 本征层。
\end{theorem}

% [Tutor Note: 定理 10.5.2 (2) 中写的是 $\Delta_{\kappa_c,y}$，而 (1) 与上下文用的是 $\Delta_{\kappa_c,I_y}$（或 $\Delta_{\kappa_c,I^0_y}$）。按前后文，这里应为 $\Delta_{\kappa_c,I_y}$；原样保留，疑为原书笔误。]

于是我们得到：若 $\chi_y=\chi|_{D^\times_y}$，则 $\gghat_{\kappa_c,y}\text{-mod}^{I_y}_{\chi_y}$ 的任一对象在函子 $\Delta_{\kappa_c,I_y}$ 下的像都属于 $\Bun_{G,y}$ 上 Hecke 本征层构成的范畴 $\mathcal{A}ut^{I_y}_{E_\chi}$。现在考虑函子 $\Delta_{\kappa_c,I^0_y}$ 到 $\gghat_{\kappa_c,y}\text{-mod}^{I^0_y}_{\chi_y}$ 上的限制。正如 §10.4.3 所讨论的，范畴 $\gghat_{\kappa_c,y}\text{-mod}^{I^0_y}_{\chi_y}$ 与相应的 $I_y$-monodromic 模范畴 $\gghat_{\kappa_c,y}\text{-mod}^{I_y,m}_{\chi_y}$ 重合。因此 $\gghat_{\kappa_c,y}\text{-mod}^{I^0_y}_{\chi_y}$ 的任一对象在函子 $\Delta_{\kappa_c,I^0_y}$ 下的像，属于 $\cD_{\kappa_c,I^0_y}\text{-mod}$ 的子范畴 $\cD^m_{\kappa_c,I^0_y}\text{-mod}$，后者的对象带有递增滤过，使相邻商都是 $\Bun_{G,y}$ 上 $\cD_{\kappa_c,I_y}$-模的拉回。这样的 $\cD_{\kappa_c,I^0_y}$-模称为 \textbf{monodromic}。

令 $\mathcal{A}ut^{I_y,m}_{E_\chi}$ 为 $\cD^m_{\kappa_c,I^0_y}\text{-mod}$ 中由以 $E_\chi$ 为本征值的 Hecke 本征层构成的子范畴。

于是我们得到函子
\begin{gather}
    \Delta_{\kappa_c,I_y}:\ \gghat_{\kappa_c,y}\text{-mod}^{I_y}_{\chi_y}
    \to \mathcal{A}ut^{I_y}_{E_\chi},
    \tag{10.5-7}\\
    \Delta_{\kappa_c,I^0_y}:\ \gghat_{\kappa_c,y}\text{-mod}^{I^0_y}_{\chi_y}
    \to \mathcal{A}ut^{I_y,m}_{E_\chi}.
    \tag{10.5-8}
\end{gather}
一个诱人的猜想是（见 [Frenkel and Gaitsgory 2006c]）：这些函子是范畴等价，至少对一般的 $\chi$ 如此。假设这成立，那么我们可以把 $\Bun_{G,I_y}$ 与 $\Bun'_{G,I^0_y}$ 上 Hecke 本征层

% --------------------------------------------------------------- p.375
的整体范畴 $\mathcal{A}ut^{I_y}_{E_\chi}$ 与 $\mathcal{A}ut^{I_y,m}_{E_\chi}$ 分别与局部范畴 $\gghat_{\kappa_c,y}\text{-mod}^{I_y}_{\chi_y}$ 及 $\gghat_{\kappa_c,y}\text{-mod}^{I^0_y}_{\chi_y}$ 等同起来。因此，我们可以用关于局部 Langlands 范畴（如 $\gghat_{\kappa_c,y}\text{-mod}^{I^0_y}_{\chi_y}$）的结果与猜想，来描述 $X$ 上带抛物结构的 $G$-丛模叠上的整体 Hecke 本征层范畴。

回忆我们对 $I^0_y$-等变模范畴的导出范畴 $D^b(\gghat_{\kappa_c,y}\text{-mod}_{\chi_y})^{I^0_y}$ 有如下猜想性描述（见公式 (10.4-14)）：
\begin{equation}
    D^b(\gghat_{\kappa_c,y}\text{-mod}_{\chi_y})^{I^0_y}
    \simeq D^b\bigl(\operatorname{QCoh}(\widetilde{\mathrm{Sp}}^{\mathrm{DG}}_{\mathrm{Res}(\chi_y)})\bigr).
    \tag{10.5-9}
\end{equation}
相应的 $I_y$-等变版本是
\begin{equation}
    D^b(\gghat_{\kappa_c,y}\text{-mod}_{\chi_y})^{I_y}
    \simeq D^b\bigl(\operatorname{QCoh}(\mathrm{Sp}^{\mathrm{DG}}_{\mathrm{Res}(\chi_y)})\bigr),
    \tag{10.5-10}
\end{equation}
其中我们把非约化的 DG Springer 纤维换成了约化的：
\[
    \cO_{\mathrm{Sp}^{\mathrm{DG}}_u}
    = \cO_{\widetilde{\cN}_{{}^L\fg}}
    \otimes_{\cO_{{}^L\fg}}\C_u .
\]
若函子 (10.5-7)、(10.5-8) 是等价，则把它们与 (10.5-9)、(10.5-10) 相结合，得到如下猜想性的范畴等价：
\begin{gather}
    D^b\bigl(\mathcal{A}ut^{I_y}_{E_\chi}\bigr)
    \simeq D^b\bigl(\operatorname{QCoh}(\mathrm{Sp}^{\mathrm{DG}}_{\mathrm{Res}(\chi_y)})\bigr),
    \tag{10.5-11}\\
    D^b\bigl(\mathcal{A}ut^{I_y,m}_{E_\chi}\bigr)
    \simeq D^b\bigl(\operatorname{QCoh}(\widetilde{\mathrm{Sp}}^{\mathrm{DG}}_{\mathrm{Res}(\chi_y)})\bigr).
    \tag{10.5-12}
\end{gather}
换言之，对应在 $y\in X$ 处驯顺分歧的局部系统的整体 Langlands 范畴（无论是否 monodromic），其导出范畴等价于它在 $y$ 处的留数所对应的 DG Springer 纤维上拟凝聚层的导出范畴（非约化或约化）。\transnote{(10.5-9) 对应 $I^0_y$ 版，取非约化 DG Springer 纤维 $\widetilde{\mathrm{Sp}}^{\mathrm{DG}}_{\mathrm{Res}(\chi_y)}$；(10.5-10) 对应 $I_y$ 版，换成约化的 $\mathrm{Sp}^{\mathrm{DG}}_{\mathrm{Res}(\chi_y)}$——p.375 用这个等式把约化纤维写成幂零锥与 $\C_u$ 的纤维积。$\mathrm{Res}(\chi_y)$ 指 oper $\chi_y$ 的留数，它给出 DG Springer 纤维所依附的幂零元。[来源：原书 p.375；(10.4-14)；§10.4.3。]}

这些等价同样应当交织范畴化仿射 Hecke 代数 $\cH(G(\cK_y),I_y)$（等价地，${}^L G$ 相应的仿射辫群）在上述范畴上的自然作用。

(10.5-11)、(10.5-12) 中出现的范畴，对 $X$ 上任一在 $y$ 处驯顺分歧的 ${}^L G$-局部系统 $E$ 其实都有意义。因此很自然地猜想这些等价一般仍然成立：
\begin{gather}
    D^b\bigl(\mathcal{A}ut^{I_y}_{E}\bigr)
    \simeq D^b\bigl(\operatorname{QCoh}(\mathrm{Sp}^{\mathrm{DG}}_{\mathrm{Res}(E)})\bigr),
    \tag{10.5-13}\\
    D^b\bigl(\mathcal{A}ut^{I_y,m}_{E}\bigr)
    \simeq D^b\bigl(\operatorname{QCoh}(\widetilde{\mathrm{Sp}}^{\mathrm{DG}}_{\mathrm{Res}(E)})\bigr).
    \tag{10.5-14}
\end{gather}
相应的局部化函子可以如下构造。假设我们可以用有限多个点 $y_1,\dots,y_n$ 的补集上的 oper $\chi$ 来表示 $X$ 上在 $y$ 处驯顺分歧的局部系统 $E$，并且 $\chi$ 到 $D^\times_{y_i}$ 的限制

% --------------------------------------------------------------- p.376
属于 $\Op^{\lambda_i}_{{}^L G}(D_{y_i})\subset\Op_{{}^L G}(D^\times_{y_i})$，其中 $\lambda_i\in P^+$。那么，与无分歧情形同样地，我们构造从 $\gghat_{\kappa_c,y}\text{-mod}^{I_y}_{\chi_y}$ 到 $\mathcal{A}ut^{I_y}_E$、以及从 $\gghat_{\kappa_c,y}\text{-mod}^{I^0_y}_{\chi_y}$ 到 $\mathcal{A}ut^{I_y,m}_E$ 的局部化函子（这里如前，$\chi_y=\chi|_{D^\times_y}$），并由此得到猜想性的等价 (10.5-13)、(10.5-14)。

这些等价可以看作\textbf{驯顺分歧局部系统的几何 Langlands 对应}。

等价 (10.5-13) 还有族版本，其中允许 $E$ 变动。它类似于局部等价的族版本 (10.4-13)。与局部情形一样，在族版本中可以避免使用 DG 概形。

上述构造可以推广到允许局部系统在有限多个点 $y_1,\dots,y_n$ 处驯顺分歧。相应的 Hecke 本征层则是 $X$ 上带抛物结构（给在 $y_1,\dots,y_n$ 处）的 $G$-丛模叠上的 D-模。这些 Hecke 本征层的非平凡例子在亏格 $0$ 时就已经出现。这些层在 [Frenkel 1995] 中被显式构造出来（另见 [Frenkel 2005a; Frenkel 2004]），它们与 Gaudin 可积系统密切相关。\transnote{亏格 $0$ 时 $X=\PP^1$，$\Bun_{G,(y_i)}$ 上的 Hecke 本征层并不平凡：这些层在 [Frenkel 1995] 中被显式构造，并与 Gaudin 可积系统密切相关——临界水平的共形块给出 Gaudin 模型的本征向量，而 Gaudin 哈密顿量作用在相应的 $D$-模上。另见 [Frenkel 2005a; Frenkel 2004] 的综述。[来源：原书 p.376；[Frenkel 1995]；[Frenkel 2005a]；[Frenkel 2004]。]}

% [Tutor Note: 「亏格 0 时的非平凡例子」指的是 $X=\PP^1$ 上带抛物结构的 $\Bun_{G,(y_i)}$ 上的 Hecke 本征层；这些层与 Gaudin 模型的关系可追溯到 Feigin--Frenkel--Reshetikhin 关于 Gaudin 模型与临界水平的工作：Gaudin 哈密顿量作用在这些层对应的 D-模上，而临界水平的共形块给出其本征向量。可参见 [Frenkel 2005a] 的综述。]

\subsection{与正则奇点的联系}

至此我们只考虑了与 $X$ 上 ${}^L G$-oper $\chi$ 相对应的 $\gghat_{\kappa_c}$-模范畴，这些 oper 除在点 $y\in X$（或若干点）之外处处正则，且其限制 $\chi|_{D^\times_y}$ 是 $\Op^{\mathrm{nilp}}_{{}^L G}(D_y)$ 中的幂零 oper。换言之，$\chi|_{D^\times_y}$ 是在 $y$ 处具有正则奇点、留数为 $\varpi(-\rho)$ 的 oper。不过，我们可以很容易地把局部化函子推广到与在 $y$ 处具有正则奇点、留数任意的 ${}^L G$-oper 相对应的 $\gghat_{\kappa_c}$-模范畴。\transnote{这里是 §10.5.3 与 §10.3--§10.4 的接合点：$D^\times_y$ 上留数为 $\varpi(-\rho)$ 的正则奇点 oper 恰好就是 §10.4 处理的幂零 oper（即 $\Op^{\mathrm{nilp}}$、$\lambda=0$ 的情形），§10.3 的局部 Langlands 范畴正是在这个情形下构造的。§10.5.3 把这套局部化函子推广到留数为 $\varpi(-\lambda-\rho)$ 的任意正则奇点 oper；其中 $\lambda=0$ 最复杂，一般 $\lambda$ 最简单。[来源：原书 p.376；§9.4.3；§10.4。]}

回忆在 §9.4.3 中，对每个具有正则奇点、留数为 $\varpi(-\lambda-\rho)$ 的 oper $\chi\in\Op^{\mathrm{RS}}_{{}^L G}(D)_{\varpi(-\lambda-\rho)}$，我们定义了以 $\chi$ 为中心特征标的 $I^0$-等变 $\gghat_{\kappa_c}$-模范畴 $\gghat_{\kappa_c}\text{-mod}^{I^0}_\chi$。$\lambda=0$ 的情形是一个“极端”情形，此时范畴 $\gghat_{\kappa_c}\text{-mod}^{I^0}_\chi$ 最为复杂。另一个“极端”是相应于一般 $\lambda$ 的一般 oper $\chi$ 的情形。此时如 §9.4.3 所见，范畴 $\gghat_{\kappa_c}\text{-mod}^{I^0}_\chi$ 相当简单：它含有由 $\fg$ 的外尔群标记的不可约对象 $\mathbb{M}_{w(\lambda+\rho)-\rho}(\chi)$，且 $\gghat_{\kappa_c}\text{-mod}^{I^0}_\chi$ 的每个对象都是这些不可约模的直和。这里 $\mathbb{M}_{w(\lambda+\rho)-\rho}(\chi)$ 是 Verma 模 $\mathbb{M}_{w(\lambda+\rho)-\rho}$（$w\in W$）对相应于 $\chi$ 的中心特征标取商所得的对象。

对其他 $\lambda$ 值，$\gghat_{\kappa_c}\text{-mod}^{I^0}_\chi$ 的结构介于这两个极端情形之间。

回忆我们有局部化函子 (10.5-6)
\[
    \Delta^\lambda_{\kappa_c,I^0_y}:\ \gghat_{\kappa_c,y}\text{-mod}^{I^0_y}
    \to \cD_{\kappa_c,I^0_y}\text{-mod}
\]
% --------------------------------------------------------------- p.377
从 $\gghat_{\kappa_c,y}\text{-mod}^{I^0_y}_{\chi_y}$ 到 $\Bun'_{G,I_y}$ 上由 $\Bun_G$ 上线丛 $K^{1/2}$ 的拉回所扭变的 D-模范畴。现在我们把这个函子限制到子范畴 $\gghat_{\kappa_c,y}\text{-mod}^{I^0_y}_{\chi_y}$，其中 $\chi_y$ 是 $D_y$ 上的 ${}^L G$-oper，在 $y$ 处有正则奇点、留数为 $\varpi(-\lambda-\rho)$。

先考虑 $\lambda\in\fh^*$ 为一般值的情形。设 $\chi_y$ 可延拓为 $X\setminus y$ 上的正则 oper $\chi$。那么与定理 10.5.2 同样地可以证明：对 $\gghat_{\kappa_c,y}\text{-mod}^{I^0_y}_{\chi_y}$ 的任一对象 $M$，相应的 $\cD_{\kappa_c,I^0_y}$-模 $\Delta_{\kappa_c,I^0_y}(M)$ 是以 $E_\chi$ 为本征值的 Hecke 本征层，这里 $E_\chi$ 是 $X$ 上以 $\chi$ 为底层的、在 $y$ 处具有正则奇点的 ${}^L G$-局部系统（若 $\chi_y$ 不能延拓到 $X\setminus y$，则如以前一样 $\Delta^\lambda_{\kappa_c,I_y}(M)=0$）。因此我们得到一个函子
\[
    \Delta_{\kappa_c,I^0_y}:\ \gghat_{\kappa_c,y}\text{-mod}^{I^0_y}_{\chi_y}
    \to \mathcal{A}ut^{I^0_y}_{E_\chi},
\]
其中 $\mathcal{A}ut^{I^0_y}_{E_\chi}$ 是 $\Bun'_{G,I_y}$ 上以 $E_\chi$ 为本征值的 Hecke 本征层范畴。

由于我们假定 oper $\chi_y$ 的留数是一般值，$E_\chi$ 围绕 $y$ 的单值性属于 ${}^L G$ 中包含 $\exp(2\pi\ii\lambda)$ 的正则半单共轭类。此时范畴 $\gghat_{\kappa_c,y}\text{-mod}^{I^0_y}_{\chi_y}$ 特别简单，如上文所讨论。我们期望函子 $\Delta_{\kappa_c,I^0_y}$ 在 $\gghat_{\kappa_c,y}\text{-mod}^{I^0_y}_{\chi_y}$ 与 $\mathcal{A}ut^{I^0_y}_{E_\chi}$ 之间建立范畴等价。\transnote{单值性的归一化极易出错：留数为 $\varpi(-\lambda-\rho)$ 的 oper 对应的局部系统，其围绕 $y$ 的单值性落在含 $\exp(2\pi\ii\lambda)$ 的正则半单共轭类中；$2\pi\ii$ 因子与 $\lambda$ 的符号、$\rho$ 平移都依赖 §9.4.3 的留数约定，不要换成其他文献中 $\exp(2\pi\ii(-\lambda-\rho))$ 的写法。[来源：原书 p.377；§9.4.3。]}

% [Tutor Note: 此处单值性的归一化极易出错：书中把留数为 $\varpi(-\lambda-\rho)$ 的 oper 与单值性属于 $\exp(2\pi\ii\lambda)$ 的共轭类相对应，因子 $2\pi\ii$ 与 $\lambda$ 的符号（含 $\rho$ 的平移）都依赖 §9.4 的约定。核对时请以 §9.4.3 的留数约定为准，不要按其他文献中 $\exp(2\pi\ii(-\lambda-\rho))$ 的写法替换。]

我们可以把这个结论表述得更漂亮。对 $M\in{}^L G$，令 $\cB_M$ 为包含 $M$ 的 Borel 子群构成的簇。注意若 $M$ 正则半单，则 $\cB_M$ 是一个点集，与 $W$ 一一对应。因此我们的猜想是：$\mathcal{A}ut^{I^0_y}_{E_\chi}$ 等价于 $\cB_M$ 上拟凝聚层的范畴 $\operatorname{QCoh}(\cB_M)$，其中 $M$ 是 $E_\chi$ 的单值性所在共轭类的代表元。\transnote{$\cB_M$ 是包含 $M$ 的 Borel 子群构成的簇，有三种形态：$M$ 正则半单时它是 $|W|$ 个点；$M$ 幂幺时它退化为 Springer 纤维 $\mathrm{Sp}_u$（$M=\exp(2\pi\ii u)$）；一般 $M$ 则是两者的“混合”。要把幂幺情形与正则半单情形统一处理，就必须把 $\cB_M$ 换成 DG 概形 $\cB^{\mathrm{DG}}_M$——这正是 p.378 最雄心猜想的内容，也与 §10.5.2 中“非约化 vs 约化 DG Springer 纤维”的区分完全平行。[来源：原书 p.377–378；§10.5.2。]}

现在考虑 $X$ 上任一在 $y\in X$ 处有正则奇点、且围绕 $y$ 的单值性为正则半单的 ${}^L G$-局部系统 $E$。于是很自然地猜想：至少当 $E$ 为一般值时，这个范畴与 $E$ 具有 oper 结构时的情形有相同的结构，即它等价于范畴 $\operatorname{QCoh}(\cB_M)$，其中 $M$ 是 $E$ 围绕 $y$ 的单值性所在共轭类的代表元。

另一方面，若围绕 $y$ 的单值性是幂幺的，则 $\cB_M$ 不是别的，正是 Springer 纤维 $\mathrm{Sp}_u$，其中 $M=\exp(2\pi\ii u)$。相应的范畴 $\mathcal{A}ut^{I^0_y}_E$ 已在 §10.5.2 中讨论过（我们期望它与 $\mathcal{A}ut^{I_y,m}_E$ 重合）。

于是我们看到，在两个“极端”情形——幂幺单值性与正则

% --------------------------------------------------------------- p.378
半单单值性——我们的猜想把 $\mathcal{A}ut^{I^0_y}_E$ 的导出范畴等同于范畴 $\operatorname{QCoh}(\cB_M)$ 的导出范畴（其中在幂幺情形，$\cB_M$ 应视为 DG 概形 $\widetilde{\mathrm{Sp}}^{\mathrm{DG}}_u$）。这就引导我们提出一个最雄心勃勃的猜想：对 $X$ 上任一在 $y\in X$ 处有正则奇点的 ${}^L G$-局部系统 $E$，$\mathcal{A}ut^{I^0_y}_E$ 的导出范畴等价于概形 $\cB_M$ 的某个合适的 DG 版本上拟凝聚层的导出范畴，其中 $M$ 是 $E$ 围绕 $y$ 的单值性所在共轭类的代表元：
\[
    D^b\bigl(\mathcal{A}ut^{I^0_y}_E\bigr)
    \simeq D^b\bigl(\operatorname{QCoh}(\cB^{\mathrm{DG}}_M)\bigr).
\]
这对多个分歧点的情形有显然的推广：此时右端取相应于各单值性的簇 $\cB^{\mathrm{DG}}_{M_i}$ 的笛卡尔积。于是，我们得到 Hecke 本征层范畴的一种猜想性实现——这些本征层的本征值是具有正则奇点的局部系统——用拟凝聚层范畴来刻画它们。

% [Tutor Note: $\cB_M$ 是 Grothendieck 意义下的「Borel 子群簇」：$\cB_M$ 正则半单时为 $|W|$ 个点，幂幺时退化为 Springer 纤维 $\mathrm{Sp}_u$（$u$ 为 $M$ 的幂零部分），一般的 $M$ 则是两者的「混合」。把 $\cB_M$ 换成 DG 概形 $\cB^{\mathrm{DG}}_M$ 是为了在幂幺情形保留非约化结构——这与 §10.5.2 中「非约化 vs 约化 DG Springer 纤维」的区分完全平行。]

下面做总结：利用临界水平的仿射 Kac--Moody 代数的表示论，我们构造了与局部 Langlands 参数（即去心圆盘上的 ${}^L G$-局部系统）相应的局部 Langlands 范畴。接着我们运用局部化函子技术，从这些局部范畴造出曲线 $X$ 上带抛物结构的 $G$-丛模叠上的整体 Hecke 本征层范畴。这些整体范畴对应整体 Langlands 参数：$X$ 上带分歧的 ${}^L G$-局部系统。我们还利用关于局部范畴结构的结果与猜想，研究了这些整体范畴。

本章我们说明了这一机制对无分歧与驯顺分歧局部系统如何运作。但我们期望，类似的方法也可用来分析具有任意分歧的局部系统。这样，仿射 Kac--Moody 代数的表示论或许有朝一日能实现那个梦想：揭开几何 Langlands 对应的奥秘。

% 本片段到此结束：§10.5 结束，第 10 章至此完结（印刷 p.378 页末）。
